% Part: computability
% Chapter: computability-theory
% Section: coding-computations

\documentclass[../../../include/open-logic-section]{subfiles}

\begin{document}

\olfileid{cmp}{thy}{cod}
\olsection{Ngode Komputasi}

Ing saben model komputasi, perkara-perkara iki bisa ditindakake. Cathetan cakupan panyunting: iki kanggo model klasik efektif standar kang dibahas ing kene. Kanggo pratelan rekursif primitif ing ngisor, dipilih pangodean program lan cathetan kang ndadekake pamriksan saben langkah lan pamacaan asil rekursif primitif. Pangodean kang mung komputabel ora otomatis menehi sipat rekursif primitif.
Perkara-perkara kasebut yaiku:
\begin{enumerate}
\item Njlentrehake \emph{definisi} fungsi komputabel
  kanthi cara sistematis. Umpamane, spesifikasi mesin
  Turing, definisi rekursif, utawa program ing basa
  pamrograman bisa dianggep menehi definisi kasebut.
\item Njlentrehake cathetan lengkap komputasi fungsi
  kang diwenehake dening definisi tartamtu kanggo
  input tartamtu. Umpamane, komputasi mesin Turing
  bisa dijlentrehake minangka rerangken konfigurasi
  (kahanan mesin lan isi pita) ing saben langkah komputasi.
\item Mriksa apa cathetan kang dianggep cathetan
  komputasi iku pancen cathetan carane fungsi komputabel
  kanthi definisi tartamtu diitung kanggo input tartamtu
  (kang ndadekake fungsi iku ditegesi, yaiku komputasine
  mandheg).
\item Njupuk nilai fungsi kanggo input tartamtu saka
  jlentrehan cathetan komputasi kang lengkap iku.
  Umpamane, isi pita ing langkah paling pungkasan
  saka komputasi mesin Turing kang mandheg iku nilaine.
\end{enumerate}

Kanthi ngode, saben jlentrehan fungsi komputabel bisa
diwenehi \emph{indeks} numerik kanthi cara kang
ndadekake instruksine bisa dipulihake saka indeks
iku sacara komputabel. Uga, cathetan lengkap sawijining
komputasi bisa dikode nganggo siji wilangan. Relasi
aritmetis kang diasilake, ``$s$~ngode cathetan komputasi
fungsi mawa indeks~$e$ kanggo input~$x$'', lan fungsi
``asil rerangken komputasi mawa kode~$s$'' banjur
komputabel; malah, kalorone rekursif primitif.

Kasunyatan dhasar iki migunani banget lan ngidini
kita mbuktekake sawetara asil kang nggumunake lan
wigati bab komputabilitas tanpa gumantung marang
model komputasi kang dipilih.

\end{document}
